\chapter{Tensorization and Low-Rank Decomposition Foundations}
\label{ch:tensors}

\section{Why tensorize EEG?}

A tensor is a multi-way array $\mathcal{X}\in\mathbb{R}^{I_1\times\cdots\times I_N}$.
The course materials distinguish order from dimension and emphasize that
flattening can obscure interactions among modes \citep{padhy2026course}. For
\EEG{}, a value can be indexed simultaneously by participant, channel,
frequency, and local time. Vectorizing these axes is always algebraically
possible, but it removes the explicit statement that neighboring frequencies,
electrodes, and time bins have different meanings.

Tensorization is therefore a hypothesis about structure. In this project:

\begin{itemize}
  \item a frequency tensor assumes subject patterns can be expressed through
  separable or multilinear channel-frequency-time factors;
  \item a local-segment logPSD tensor preserves signed deviations from a
  training mean and supports orthogonal TT/Tucker projection;
  \item a waveform-delay tensor assumes short-lag covariance contains useful
  information even when absolute waveform phase is not aligned across people.
\end{itemize}

\section{Rank-one structure and CP decomposition}

A third-order rank-one tensor is the outer product
$\bm{a}\circ\bm{b}\circ\bm{c}$ with entries
$x_{ijk}=a_i b_j c_k$. CP represents an order-$N$ tensor as a sum of $R$
rank-one terms \citep{kolda2009tensor}:

\begin{equation}
\mathcal{X}\approx\sum_{r=1}^{R}
\bm{a}^{(1)}_r\circ\bm{a}^{(2)}_r\circ\cdots\circ\bm{a}^{(N)}_r.
\label{eq:cp}
\end{equation}

For the V1 subject-channel-frequency-time tensor, the subject factor can be
used as a compact embedding. Non-negativity improves parts-based interpretability
for non-negative power data, but it is an optimization constraint, not a
guarantee of identifiability or classification performance.

The preliminary graph-regularized V1 objective has the conceptual form

\begin{equation}
\min_{A,B,C,D\geq0}
\frac{1}{2}\left\|\mathcal{X}-[[A,B,C,D]]\right\|_F^2
+\lambda\operatorname{tr}(B^T L B),
\label{eq:graphcp}
\end{equation}

where $B$ is the channel factor and $L$ is an electrode-graph Laplacian. The
implemented multiplicative update mixes reconstruction and graph penalties,
but its recorded objective scales those terms inconsistently under factor
renormalization. The corresponding result is retained as preliminary rather
than presented as validated graph regularization.

\section{Tucker decomposition and HOSVD}

Tucker decomposition uses a core tensor and one factor matrix per mode:

\begin{equation}
\mathcal{X}\approx\mathcal{G}\times_1 U^{(1)}\times_2 U^{(2)}
\cdots\times_N U^{(N)}.
\label{eq:tucker}
\end{equation}

Unlike CP's single component count, Tucker can choose a different multilinear
rank for each mode. V2 implements a training-only HOSVD-like projection: each
factor is formed from leading eigenvectors of a mode covariance, and their
Kronecker-structured combination defines a reusable orthonormal basis. V3 uses
channel and lag factors to form waveform-energy filters. Tucker is a natural
course bridge from matrix PCA to structured higher-order compression
\citep{kolda2009tensor,padhy2026lowrank}.

\section{Tensor Train decomposition}

Tensor Train writes each entry as a product of small cores:

\begin{equation}
x_{i_1\ldots i_N}=G_1[i_1]G_2[i_2]\cdots G_N[i_N],
\label{eq:tt}
\end{equation}

with boundary ranks equal to one. TT-SVD obtains cores from successive low-rank
unfoldings and controls storage through intermediate TT ranks
\citep{oseledets2011tt}. In V2, feature modes are placed before observations;
left-orthogonal cores define a common training-derived basis. This avoids
comparing separately factorized subjects whose cores may differ by arbitrary
signs, rotations, or gauges.

For the V2 $19\times59\times3$ window tensor, the basis is assembled as

\begin{equation}
B_{(i,j,k),d}=\sum_{a,b}G_1(1,i,a)G_2(a,j,b)G_3(b,k,d).
\end{equation}

A centered held-out tensor is projected by matrix multiplication with $B$; the
basis is never refit on held-out people.

\section{Welch and spectrogram tensorization}

Welch's method averages modified periodograms from windowed, overlapping
segments to reduce spectral-estimate variance \citep{welch1967fft}. The project
uses a 500-sample Hann segment at 250 Hz, 250-sample overlap, and 500-point FFT,
giving 0.5 Hz bins. Frequencies from 1 to 30 Hz yield 59 bins.

The V1 relative-power map for subject $n$, channel $c$, frequency $f$, and time
bin $q$ can be written as

\begin{equation}
X_{ncfq}=\operatorname{median}_{w\in q}
\sqrt{\frac{P_{ncfw}}{\sum_{c',f'}P_{nc'f'w}+\epsilon}}.
\end{equation}

V2 instead preserves signed log absolute power:

\begin{equation}
Z_{w,c,f,s}=\log_{10}\left(\max(P_{w,c,f,s},10^{-12})\right),
\end{equation}

where $s\in\{1,2,3\}$ indexes the three local overlapping two-second segments
within a four-second window. This representation can express increases and
decreases after centering, so non-negativity is neither required nor desirable.

\section{Waveform-delay tensor and exact moment compression}

Each clean four-second window is anti-alias resampled from 250 to 125 Hz,
producing 500 samples. A 16-lag sliding view forms
$\mathcal{D}\in\mathbb{R}^{19\times16\times485}$, where lag spacing is 8 ms and
the first-to-last span is 120 ms. Each channel-lag row is centered over local
positions.

Let $z_p\in\mathbb{R}^{304}$ be the vectorized channel-lag sample at local
position $p$. The per-window moment is

\begin{equation}
S_w=\frac{1}{485}\sum_{p=1}^{485} z_p z_p^T,
\end{equation}

and the subject moment is the equal-window average. For any filter
$b\in\mathbb{R}^{304}$, its mean squared response is exactly

\begin{equation}
E(b)=\frac{b^T S b}{b^T b}.
\label{eq:moment-energy}
\end{equation}

Thus the stored $88\times304\times304$ moments are exact sufficient statistics
for the pooled squared response of any linear delay filter. They avoid storing
all delay tensors, but they deliberately discard local-position and long-range
sequence order. The waveform branch is therefore a pooled second-order model,
not a general raw-EEG neural sequence model.

\begin{figure}[!htbp]
\centering
\begin{tikzpicture}[node distance=6mm and 13mm,
block/.append style={text width=65mm,minimum height=13mm},
process/.append style={text width=65mm,minimum height=13mm}]
\node[block,text width=94mm] (wave) at (0,0) {One four-second EEG window: $19\times1000$ at 250 Hz};
\node[process] (psd) at (-3.9,-1.8) {Spectral path\\Welch / spectrogram: 1--30 Hz, 0.5 Hz steps};
\node[process] (delay) at (3.9,-1.8) {Waveform-delay path\\Resample to 125 Hz; form 16-lag views};
\node[block,below=of psd] (spec) {V1 cohort: $69\times19\times59\times10$\\V2 window: $19\times59\times3$};
\node[block,below=of delay] (dt) {Delay window: $19\times16\times485$\\Exact moment: $304\times304$};
\node[process,below=of spec] (lowrank) {PCA / NMF / CP / Tucker / TT\\Training-only basis and subject pooling};
\node[process,below=of dt] (energy) {TT / Tucker filters\\Pooled log-energy features};
\node[block,text width=94mm] (features) at (0,-7.7) {One feature vector per participant\\Training-fitted scaling / selection, then classifier};
\draw[arrow] (wave.south) -- ++(0,-0.12) -| (psd.north);
\draw[arrow] (wave.south) -- ++(0,-0.12) -| (delay.north);
\draw[arrow] (psd)--(spec);
\draw[arrow] (delay)--(dt);
\draw[arrow] (spec)--(lowrank);
\draw[arrow] (dt)--(energy);
\draw[arrow] (lowrank.south) -- ++(0,-0.4) -| (features.north);
\draw[arrow] (energy.south) -- ++(0,-0.4) -| (features.north);
\end{tikzpicture}
\caption{Implemented spectral and waveform-delay tensorizations, with their actual array dimensions. The two branches are alternative representations; fusion is evaluated separately in later experiments.}
\label{fig:tensorization}
\end{figure}


